\documentclass[a4paper,11pt]{article}
\usepackage[utf8]{inputenc}
\usepackage[T1]{fontenc}
\usepackage[ngerman]{babel}
\usepackage{geometry}
\geometry{left=2.5cm, right=2.5cm, top=2.5cm, bottom=2.5cm}
\usepackage{longtable}
\usepackage{booktabs}
\usepackage{array}
\usepackage{xcolor}
\usepackage{colortbl}
\usepackage{hyperref}
\usepackage{fancyhdr}
\usepackage{listings}
\usepackage{tabularx}
\usepackage{enumitem}
\usepackage{graphicx}

\definecolor{headerblue}{HTML}{1565C0}
\definecolor{tableheader}{HTML}{E3F2FD}
\definecolor{codebg}{HTML}{F5F5F5}
\definecolor{codegreen}{HTML}{2E7D32}
\definecolor{codegray}{HTML}{757575}

\lstset{
  backgroundcolor=\color{codebg},
  basicstyle=\ttfamily\small,
  keywordstyle=\color{headerblue}\bfseries,
  commentstyle=\color{codegreen},
  stringstyle=\color{codegray},
  breaklines=true,
  frame=single,
  rulecolor=\color{codegray!30},
  numbers=none,
  showstringspaces=false,
  tabsize=2,
  language=C++
}

\hypersetup{colorlinks=true,linkcolor=headerblue,urlcolor=headerblue}

\pagestyle{fancy}
\fancyhf{}
\fancyhead[L]{\textbf{INSPECTAIR}}
\fancyhead[R]{Design Description v2.0}
\fancyfoot[C]{\thepage}
\renewcommand{\headrulewidth}{0.4pt}

\setlength{\parindent}{0pt}
\setlength{\parskip}{6pt}

\begin{document}

\begin{center}
{\LARGE\textbf{INSPECTAIR}}\\[0.3cm]
{\Large Design Description}\\[0.3cm]
{\large System- und Software-Architektur}\\[0.5cm]
\begin{tabular}{ll}
\textbf{Projekt:} & InspectAir -- Luftqualitätsmessgerät \\
\textbf{Version:} & 2.0 \\
\textbf{Datum:} & Februar 2026 \\
\textbf{Autoren:} & Maximilian Liebl, Cedric Stroh, Jannis Messer \\
\textbf{Modul:} & Mikrocomputertechnik 1, DHBW Ravensburg \\
\textbf{Dozent:} & Thomas Stingl (Airbus) \\
\end{tabular}
\end{center}

\vspace{0.5cm}
\noindent\textit{Dieses Dokument wurde unter Verwendung von KI-Tools (Claude, Anthropic) zur Überprüfung erstellt.}
\vspace{0.3cm}

\tableofcontents
\newpage

% ============================================================
\section{Projektbeschreibung}

\subsection{Überblick}
InspectAir ist ein Luftqualitätsmonitor für Innenräume, basierend auf dem ESP32-S3. Das System erfasst kontinuierlich sechs Umweltparameter (CO\textsubscript{2}, PM2.5, PM10, VOC, Temperatur, Luftfeuchtigkeit), erkennt Personen per 24\,GHz-Radar und stellt alle Daten auf einem 4-Zoll-TFT-Display dar. Durch eine vierstufige Farbcodierung nach WHO-2021- und UBA-Grenzwerten können kritische Werte auf einen Blick erkannt werden.

\subsection{Kernfunktionen}
Das System bietet folgende Hauptfunktionen:
\begin{itemize}[nosep]
\item Messung von CO\textsubscript{2}, Feinstaub (PM1.0/PM2.5/PM10), VOC, Temperatur und Luftfeuchtigkeit
\item Echtzeit-Anzeige auf 4-Zoll-Display (480$\times$320 Pixel) über LVGL 9.2
\item Fünf wählbare Bildschirmansichten (Overview, Detail, Analog, Bubble, Tree-Animation)
\item Vierstufige Farbcodierung (Grün/Gelb/Orange/Rot) nach WHO/UBA-Grenzwerten
\item Präsenzerkennung via 24\,GHz FMCW-Radar für automatisches Display-Management
\item Energiesparmodi: Display-Dimming, Display-Off und Light Sleep
\item WiFi-Anbindung für NTP-Zeitsynchronisation und optionale Web-Fernsteuerung
\end{itemize}

% ============================================================
\section{Systemarchitektur}

\subsection{Hardware-Blockdiagramm}

Das System besteht aus einem ESP32-S3 Mikrocontroller als zentraler Steuereinheit, verbunden mit fünf Sensoren über drei verschiedene Busse sowie einem TFT-Display über SPI:

\begin{verbatim}
  ┌─────────────┐     SPI      ┌─────────────────┐
  │  ST7796S    │◄────────────►│                 │
  │  480x320    │              │   ESP32-S3      │
  │  (LovyanGFX)│              │   N8R8 DevKit   │
  └─────────────┘              │   (Waveshare)   │
                               │                 │
  ┌─────────────┐     I2C      │  Flash: 8MB QIO │
  │ AHT20       │◄────────────►│  PSRAM: 8MB OPI │
  │ SGP40       │  (4.7kΩ PU)  │                 │
  └─────────────┘              │                 │
                               │                 │
  ┌─────────────┐   UART1      │                 │
  │ PMS5003     │◄────────────►│                 │
  └─────────────┘   9600 Bd    │                 │
                               │                 │
  ┌─────────────┐  SoftSerial  │                 │
  │ MH-Z19C    │◄────────────►│                 │
  │ (+220µF)    │   9600 Bd    │                 │
  └─────────────┘              │                 │
                               │                 │
  ┌─────────────┐   UART2      │                 │
  │ LD2410C     │◄──[TXS0108E]►│                 │
  │ (+220µF)    │  256000 Bd   │                 │
  └─────────────┘              └─────────────────┘
\end{verbatim}

\subsection{Bus-Zuordnung}

\begin{tabularx}{\textwidth}{llXl}
\toprule
\rowcolor{tableheader}
\textbf{Bus} & \textbf{Sensor} & \textbf{Pins} & \textbf{Baudrate} \\
\midrule
SPI & ST7796S Display & MOSI:11, MISO:13, SCLK:12, CS:14, DC:9, RST:46 & -- \\
I2C & AHT20 + SGP40 & SDA:8, SCL:18 (4{,}7\,k$\Omega$ Pull-ups) & -- \\
UART1 & PMS5003 & RX:16, TX:17 & 9600 \\
SoftwareSerial & MH-Z19C & RX:4, TX:5 & 9600 \\
UART2 & LD2410C & RX:7, TX:6 (via Level Shifter) & 256000 \\
GPIO & UI-Button & GPIO1 (Active LOW, interner Pull-up) & -- \\
GPIO & Backlight & GPIO3 (via LovyanGFX PWM) & -- \\
\bottomrule
\end{tabularx}

% ============================================================
\section{Software-Architektur}

\subsection{Schichtenmodell}

Die Software folgt einem 3-Schichten-Modell:

\begin{tabularx}{\textwidth}{lX}
\toprule
\rowcolor{tableheader}
\textbf{Schicht} & \textbf{Module} \\
\midrule
\textbf{Application Layer} & \texttt{main.cpp} -- Hauptschleife, Timing, Button-Handler, Watchdog \\
\textbf{Service Layer} & \texttt{ui\_manager} (5 Screens), \texttt{sensor\_filter} (Moving Average), \texttt{sensor\_history} (24h-Daten), \texttt{power\_manager} (Energieverwaltung), \texttt{WifiClock} (NTP), \texttt{web\_remote} (HTTP-Server) \\
\textbf{Driver Layer} & \texttt{lvgl\_driver} (LVGL~9 + LovyanGFX), \texttt{aht\_sgp} (I2C), \texttt{mhz19c} (SoftSerial), \texttt{pms5003} (UART1), \texttt{ld2410c} (UART2) \\
\bottomrule
\end{tabularx}

\subsection{Dateistruktur}

\begin{tabularx}{\textwidth}{lX}
\toprule
\rowcolor{tableheader}
\textbf{Datei} & \textbf{Funktion} \\
\midrule
\texttt{main.cpp} & Hauptprogramm: setup(), loop(), Button-Handler, Sensor-Timing \\
\texttt{pins.h} & Zentrale GPIO-Pin-Definitionen \\
\texttt{sensor\_types.h} & Datenstrukturen: \texttt{SensorReadings}, \texttt{AHT20\_Data}, etc. \\
\texttt{colors.h} & RGB565-Farben, Grenzwerte, Klassifizierungsfunktionen \\
\texttt{display\_config.h} & LovyanGFX-Konfiguration für ST7796S \\
\texttt{app\_config.h} & Applikations-Konfiguration \\
\texttt{lvgl\_driver.cpp/h} & LVGL~9 Display-Treiber, Flush-Callback, Brightness-Steuerung \\
\texttt{ui\_manager.cpp/h} & Multi-Screen UI System (5 Screens), Update-Funktionen \\
\texttt{aht\_sgp.cpp/h} & AHT20 + SGP40 I2C-Treiber \\
\texttt{mhz19c.cpp/h} & MH-Z19C CO\textsubscript{2}-Treiber über SoftwareSerial \\
\texttt{pms5003.cpp/h} & PMS5003 Feinstaub-Treiber über UART1 \\
\texttt{ld2410c.cpp/h} & LD2410C Radar-Treiber über UART2 \\
\texttt{sensor\_filter.cpp/h} & Moving-Average-Filter für Sensordaten \\
\texttt{sensor\_history.cpp/h} & 24-Stunden-Datenhistorie \\
\texttt{power\_manager.cpp/h} & Display-Zustandsmaschine, Fade-Effekte, Light Sleep \\
\texttt{WifiClock.cpp/h} & WiFi-Verbindung und NTP-Zeitsynchronisation \\
\texttt{web\_remote.cpp/h} & HTTP-Server für Fernsteuerung per Smartphone \\
\bottomrule
\end{tabularx}

\subsection{Zentrale Datenstruktur}

Alle Sensordaten werden in einer aggregierten Struktur zusammengefasst (definiert in \texttt{sensor\_types.h}):

\begin{lstlisting}
typedef struct {
  AHT20_Data aht;       // Temperatur/Feuchte
  SGP40_Data sgp;       // VOC-Index + Rohwert
  PMS5003_Data pms;     // PM1.0, PM2.5, PM10 (AE + SP)
  MHZ19C_Data mhz;      // CO2 in ppm + valid-Flag
  LD2410C_Data radar;   // Praesenz, Entfernung, Bewegung
  uint32_t timestamp;   // Zeitstempel (millis)
} SensorReadings;
\end{lstlisting}

\subsection{Display-Architektur}

Die Display-Ausgabe basiert auf einer zweistufigen Architektur:

\textbf{Schicht 1 -- LovyanGFX} (\texttt{display\_config.h}, \texttt{lvgl\_driver.cpp}): Konfiguriert die SPI-Kommunikation mit dem ST7796S-Controller, verwaltet den Backlight-PWM über \texttt{setBrightness(0-255)} und stellt den Flush-Callback bereit. Dual-Buffer-Strategie mit PSRAM-Allokation (\texttt{LVGL\_BUF\_SIZE = 480*32}).

\textbf{Schicht 2 -- LVGL 9.2} (\texttt{ui\_manager.cpp/h}): Verwaltet fünf verschiedene Bildschirme:

\begin{tabularx}{\textwidth}{llX}
\toprule
\rowcolor{tableheader}
\textbf{Enum} & \textbf{Screen} & \textbf{Beschreibung} \\
\midrule
\texttt{UI\_SCREEN\_TREE} & Baum-Animation & Startbildschirm mit animierter Baumgrafik \\
\texttt{UI\_SCREEN\_OVERVIEW} & Übersicht & Große AQI-Anzeige rechts, 2 Kacheln (Temp, Feuchte) links \\
\texttt{UI\_SCREEN\_DETAIL} & Detail & Kleine AQI, 4 Kacheln mit allen Messwerten \\
\texttt{UI\_SCREEN\_ANALOG} & Analog & Analoge Cockpit-Instrumente \\
\texttt{UI\_SCREEN\_BUBBLE} & Bubbles & Dynamische Kreise (Modern Bubbles Design) \\
\bottomrule
\end{tabularx}

Custom Fonts mit deutschen Umlauten sind in fünf Größen integriert: 12, 16, 20, 28 und 48\,pt.

\subsection{Zustandsmaschine -- Power Management}

Der \texttt{PowerManager} implementiert eine 4-Zustands-Maschine:

\begin{verbatim}
  ACTIVE ──(45s ohne Präsenz)──► DIMMED ──(5 Min)──► OFF ──► SLEEPING
    ▲                              │                   │        │
    └──────────────────────────────┴───────────────────┴────────┘
                    (Präsenz erkannt → wakeUp)
\end{verbatim}

\begin{tabularx}{\textwidth}{lXl}
\toprule
\rowcolor{tableheader}
\textbf{Zustand} & \textbf{Beschreibung} & \textbf{Helligkeit} \\
\midrule
\texttt{STATE\_ACTIVE} & Normalbetrieb, volle Displayhelligkeit & 255 (100\,\%) \\
\texttt{STATE\_DIMMED} & Nach 45\,s ohne Präsenz, reduzierte Helligkeit & 30 ($\sim$12\,\%) \\
\texttt{STATE\_OFF} & Nach 5\,Min ohne Präsenz, Display ausgeschaltet & 0 \\
\texttt{STATE\_SLEEPING} & ESP32 Light Sleep, minimaler Stromverbrauch & 0 \\
\bottomrule
\end{tabularx}

Übergänge zwischen den Helligkeitsstufen erfolgen als Non-Blocking-Fade über \texttt{\_startFade()} und \texttt{\_updateFade()}.

% ============================================================
\section{Bibliotheken und Abhängigkeiten}

\begin{tabularx}{\textwidth}{llX}
\toprule
\rowcolor{tableheader}
\textbf{Bibliothek} & \textbf{Version} & \textbf{Verwendung} \\
\midrule
LovyanGFX & 1.1.16 & Display-Treiber für ST7796S, Backlight-PWM \\
LVGL & 9.2.2 & UI-Framework, Widgets (Arc, Bar, Button, Image, Label, Flex) \\
Adafruit AHTX0 & latest & AHT20 Temperatur-/Feuchte-Sensor \\
Adafruit SGP40 & latest & SGP40 VOC-Sensor mit VOC-Index-Algorithmus \\
MH-Z19 & 1.5.4 & MH-Z19C CO\textsubscript{2}-Sensor-Treiber \\
PMS Library & 1.1.0 & PMS5003 Feinstaub-Sensor-Treiber \\
ld2410 & 0.1.3 & LD2410C Radar-Sensor-Treiber \\
EspSoftwareSerial & 8.2.0 & Software-UART für MH-Z19C \\
\bottomrule
\end{tabularx}

% ============================================================
\section{Kritische Designentscheidungen}

\begin{tabularx}{\textwidth}{lX}
\toprule
\rowcolor{tableheader}
\textbf{Entscheidung} & \textbf{Begründung} \\
\midrule
220\,\textmu F Stützkondensatoren & Puffern Stromspitzen der 5V-Sensoren (MH-Z19C, LD2410C), die sonst Systemabstürze verursachen. \\
Level Shifter (TXS0108E) & LD2410C sendet 5V UART-Signale, die den 3{,}3V-ESP32 ohne Level Shifter zerstören würden. \\
SoftwareSerial für MH-Z19C & UART0 ist USB-Monitor, UART1 ist PMS5003, UART2 ist LD2410C -- daher Software-UART (9600\,Baud ist unkritisch). \\
LVGL 9.2 statt 8.3 & Modernere API, besseres Speichermanagement, native Custom-Font-Unterstützung mit Umlauten. \\
PSRAM für LVGL-Buffer & Dual-Buffer in PSRAM entlastet den internen SRAM für Anwendungslogik. \\
Light Sleep statt Deep Sleep & Schnellere Aufwachzeit ($\sim$2\,ms vs. $\sim$300\,ms), RAM bleibt erhalten, Radardaten persistent. \\
PETG 3D-Druck-Gehäuse & Temperaturbeständiger als PLA, bessere Schlagfestigkeit, gute Druckbarkeit mit Bambu Studio. \\
\bottomrule
\end{tabularx}

% ============================================================
\section{Änderungshistorie}

\begin{tabularx}{\textwidth}{l l l X}
\toprule
\rowcolor{tableheader}
\textbf{Version} & \textbf{Datum} & \textbf{Autor} & \textbf{Änderung} \\
\midrule
1.0 & Januar 2026 & Team & Erstversion \\
2.0 & Februar 2026 & Team & LVGL 8.3$\to$9.2, LovyanGFX 1.1.16, 5 Screens, SoftwareSerial, PETG-Gehäuse, aktualisierte Pin-Belegung, 4-Zustands-PowerManager, Web Remote. \\
\bottomrule
\end{tabularx}

\end{document}
